\section{Conclusion}\label{sec:conclusion}

This paper studies standards coalitions when governments inside a non-singleton coalition choose additional harmonization depth and firms can subsequently reposition products. The central mechanism is strategic re-differentiation. A regional standards union compresses standard-based differentiation between its members, but those member firms can respond by moving apart on another horizontal product dimension. When redesign costs are in the relevant regular range, this response can raise domestic producer rents enough to reverse the member-country welfare ranking.

The result is deliberately benchmarked against models that endogenize only one margin. With endogenous non-singleton harmonization depth but fixed product positions, international standardization is stable at the canonical regular region. With endogenous product positions but zero additional harmonization depth, international standardization is again stable. Only the full interaction generates a regional-union outcome. This is the paper's main theoretical contribution.

The welfare analysis shows why the political and global rankings can diverge. The member-country gain is driven by domestic producer rents that outweigh a consumer-surplus loss, while price transfers cancel in world welfare. At the canonical witness, the regional union is politically stable even though international standardization yields higher global welfare. Firms also re-differentiate more than a constrained social planner after regional integration is fixed.

These results do not imply that regional standardization is always inefficient or that all interoperability policies induce re-differentiation. The headline reversal is a constructive result on a verified regular region, and whole-circle continuation validity is established computationally at the canonical primitives over the repaired feasible policy domain rather than by a global analytic theorem. The broader implication is more specific: when policy standardizes one interface dimension but leaves firms discretion elsewhere, strategic product redesign can alter the coalition incentives that motivated the policy comparison in the first place. Standards policy and product differentiation should therefore be analyzed jointly when the firms' post-standardization design response is economically important.